% PAGES: 47-52
% Page 52 (book p.36) is blank in the source -- transitional blank verso
% before Chapter 2, which begins on the next page. Nothing to transcribe
% on that page.

\section{Exercises}

\begin{exercises}

\item Let
\[
A \;=\; \begin{pmatrix} 1 & -1 & 2 & 5 & 4 \\ 3 & -2 & 1 & 4 & 2 \\ 0 & 1 & 2 & -1 & 3 \\ -5 & 4 & 2 & -4 & 3 \end{pmatrix}.
\]
Show that the column rank and the row rank of $A$ are both equal to $3$.

\item Let $x$ and $y$ be column vectors of size $n$, and let $I$ be the
identity matrix of size $n$.
\begin{subparts}
\item If $y^{t}x \ne -1$, show that
\[
(I+xy^{t})^{-1} \;=\; I \;-\; \frac{1}{1+y^{t}x}\,xy^{t}.
\]
In other words, show that
\[
\left(I - \frac{1}{1+y^{t}x}\,xy^{t}\right)(I+xy^{t}) \;=\; I.
\]
\item Show that the matrix $I+xy^{t}$ is nonsingular if and only if
$y^{t}x\ne -1$.
\end{subparts}

\item
\begin{subparts}
\item Use induction to show that
\[
\left(\prod_{i=1}^n A_i\right)^{t} \;=\; \prod_{i=1}^n A_{n+1-i}^{t}
\]
for any $m_i\times n_i$ matrices $A_i$, $i=1:n$, with $n_i=m_{i+1}$ for
$i=1:(n-1)$.
\item Show that
\[
\left(\prod_{i=1}^n A_i\right)^{-1} \;=\; \prod_{i=1}^n A_{n+1-i}^{-1}
\]
for any nonsingular square matrices $A_i$ of the same size.
\end{subparts}

\item Let $D=\diag(d_i)_{i=1:n}$ be a diagonal matrix of size $n$ with
distinct diagonal entries, i.e., such that $d_j\ne d_k$, for any
$1\le j\ne k\le n$. If $A$ is a square matrix of size $n$, show that
$AD=DA$ if and only if the matrix $A$ is diagonal.

\item Use the fact that $D_1D_2=D_2D_1$ for any two diagonal matrices
$D_1$ and $D_2$ of the same size to show that
\[
\prod_{i=1}^n D_i \;=\; \prod_{i=1}^n D_{p(i)},
\]
for any one-to-one function $p:\{1,2,\ldots n\}\to\{1,2,\ldots n\}$,
where $D_i$, $i=1:n$, are diagonal matrices of the same size.

\item
\begin{subparts}
\item Let $A$ be an $n\times n$ matrix and let $L$ be an $n\times n$
nonsingular lower triangular matrix. Show that, if $LA$ is a lower
triangular matrix, then $A$ is lower triangular. Show that, if $AL$ is a
lower triangular matrix, then $A$ is lower triangular.
\item Let $A$ be an $n\times n$ matrix and let $U$ be an $n\times n$
nonsingular upper triangular matrix. Show that, if $UA$ is an upper
triangular matrix, then $A$ is upper triangular. Show that, if $AU$ is an
upper triangular matrix, then $A$ is upper triangular.
\end{subparts}

\item Let $A$ be a nonsingular matrix, and let $k$ be a positive integer.
Define $A^{-k}$ as the $k$--th power of the inverse matrix of $A$, i.e.,
let $A^{-k}=\left(A^{-1}\right)^{k}$. Show that this definition is
consistent, i.e., show that
\[
A^{k}\cdot A^{-k} \;=\; A^{-k}\cdot A^{k} \;=\; I.
\]

\item
\begin{subparts}
\item Let
\[
M \;=\; \begin{pmatrix} 0&0&0&0\\ 3&0&0&0\\ 1&-1&0&0\\ -1&2&1&0 \end{pmatrix}.
\]
Compute $M^2$, $M^3$, $M^4$.
\item Let
\[
C \;=\; I+M \;=\; \begin{pmatrix} 1&0&0&0\\ 3&1&0&0\\ 1&-1&1&0\\ -1&2&1&1 \end{pmatrix}.
\]
Compute $C^m$, where $m\ge 2$ is a positive integer.
\end{subparts}
\noindent Hint: Recall that, if $A$ and $B$ are square matrices
of the same size such that $AB=BA$, then the following version of the
binomial formula holds true:
\begin{equation}
(A+B)^m \;=\; \sum_{j=0}^m \binom{m}{j} A^{j}B^{m-j},
\end{equation}
where $m$ is a positive integer and the binomial coefficient
$\binom{m}{j}$ is given by
\[
\binom{m}{j} \;=\; \frac{m!}{j!\,(m-j)!},
\]
where $k!=1\cdot 2\cdot\ldots\cdot k$. Also, note that $A^0=B^0=I$.

\item Let $L$ be an $n\times n$ lower triangular matrix with entries equal
to $0$ on the main diagonal, i.e., with $L(i,i)=0$ for $i=1:n$.
\begin{subparts}
\item Show that $L^n = 0$;
\item Compute $(I+L)^m$ in terms of $L, L^2, \ldots, L^{n-1}$, where
$m\ge n$ is a positive integer.
\end{subparts}
\noindent Hint: Use the binomial formula (1.115).

\item Let $A$ and $B$ be square matrices of the same size with
nonnegative entries and such that the sum of the entries in each row is
equal to $1$. Show that the matrix $AB$ has the same properties, i.e.,
show that all the entries of the matrix $AB$ are nonnegative and the sum
of the entries in each row of $AB$ is equal to $1$.

\noindent Note: A matrix with nonnegative entries such that the
sum of the entries in each row is equal to $1$ is called a probability
matrix.

\item The covariance matrix of five random variables is
\[
\Sigma \;=\; \begin{pmatrix}
1 & -0.525 & 1.375 & -0.075 & -0.75 \\
-0.525 & 2.25 & 0.1875 & 0.1875 & -0.675 \\
1.375 & 0.1875 & 6.25 & 0.4375 & -1.875 \\
-0.075 & 0.1875 & 0.4375 & 0.25 & 0.3 \\
-0.75 & -0.675 & -1.875 & 0.3 & 9
\end{pmatrix}.
\]
Find the correlation matrix of these random variables.

\item The correlation matrix of five random variables is
\[
\Omega \;=\; \begin{pmatrix}
1 & -0.25 & 0.15 & -0.05 & -0.30 \\
-0.25 & 1 & -0.10 & -0.25 & 0.10 \\
0.15 & -0.10 & 1 & 0.20 & 0.05 \\
-0.05 & -0.25 & 0.20 & 1 & 0.10 \\
-0.30 & 0.10 & 0.05 & 0.10 & 1
\end{pmatrix}.
\]
\begin{subparts}
\item Compute the covariance matrix of these random variables if their
standard deviations are $0.25$, $0.5$, $1$, $2$, and $4$, in this order.
\item Compute the covariance matrix of these random variables if their
standard deviations are $4$, $2$, $1$, $0.5$, and $0.25$, in this order.
\end{subparts}

\item The file \textit{indeces-jul26-aug9-2012.xlsx} from
fepress.org/nla-primer contains the July 26, 2012 -- August 9, 2012 end of
day values of Dow Jones, Nasdaq, and S\&P 500.
\begin{subparts}
\item Compute the daily percentage returns of the three indices over the
given time period.
\item Compute the covariance matrix of the daily percentage returns of the
three indices.
\item Compute the daily log returns of the three indices over the given
time period.
\item Compute the covariance matrix of the daily log returns of the three
indices.
\end{subparts}

\noindent Note: The percentage return and the log return between
times $t_1$ and $t_2$ of an asset with price $S(t)$ at time $t$ are given
by
\[
\frac{S(t_2)-S(t_1)}{S(t_1)} \qquad \text{and} \qquad
\ln\!\left(\frac{S(t_2)}{S(t_1)}\right),
\]
respectively.

\item The file \textit{indices-july2011.xlsx} from fepress.org/nla-primer
contains the January 2011 -- July 2011 end of day values of nine major US
indeces.
\begin{subparts}
\item Compute the sample covariance matrix of the daily percentage
returns of the indeces, and the corresponding sample corelation matrix.

Compute the sample covariance and correlation matrices for daily log
returns, and compare them with the corresponding matrices for daily
percentage returns.
\item Compute the sample covariance matrix of the weekly percentage
returns of the indeces, and the corresponding sample corelation matrix.

Compute the sample covariance and correlation matrices for weekly log
returns, and compare them with the corresponding matrices for weekly
percentage returns.
\item Compute the sample covariance matrix of the monthly percentage
returns of the indeces, and the corresponding sample corelation matrix.

Compute the sample covariance and correlation matrices for monthly log
returns, and compare them with the corresponding matrices for monthly
percentage returns.
\item Comment on the differences between the sample covariance and
correlation matrices for daily, weekly, and monthly returns.
\end{subparts}

\item In three months, the value of an asset with spot price \$50 will be
either \$60 or \$45. The continuously compounded risk--free rate is
$6\%$. Consider the one period market model with two securities, i.e.,
cash and the asset, and two states, i.e., asset value equal to \$60 and
asset value equal to \$45, in three months.
\begin{subparts}
\item Find the payoff matrix of this model.
\item Is this one period market complete, i.e., is the payoff matrix
nonsingular?
\item How do you replicate a three months at-the-money put option on
this asset, using the cash and the underlying asset?
\end{subparts}

\item In six months, the price of an asset with spot price \$40 will be
either \$30, \$35, \$40, \$42, \$45, or \$50. Consider a one period
market model with six states in six months corresponding to the six
possible values of the asset in six months, and with the following four
securities:
\begin{itemize}
\item cash;
\item asset;
\item six months at-the-money call option with strike \$40 on the
asset;
\item six months at-the-money put option with strike \$40 on the asset.
\end{itemize}
The continuously compounded risk--free interest rate is constant and
equal to $6\%$.
\begin{subparts}
\item Find the payoff matrix of this model.
\item Is this one period market model complete?
\item Are the four securities non--redundant?
\end{subparts}

\end{exercises}
